\documentclass[12pt]{article}
\usepackage[utf8]{inputenc}
\usepackage[T1]{fontenc}
\usepackage{graphicx}
\usepackage{amsmath}
\usepackage{booktabs}
\usepackage{caption}
\usepackage{float,pdfescape}

\title{\\Análise de Otimização em Processamento Alto Desempenho}

\begin{document}

\maketitle

% PART I - Cell 0 (markdown)
\section*{Parte I -- Análise conceitual}

\subsection*{Questão 1 -- Otimização não significa apenas "fazer menos contas"}

Supondo que A e B rodam no mesmo computador. O programa B pode ter uma melhor localidade espacial que o A, como é o exemplo de casos de acesso de matrizes por colunas(pior localidade) e por linhas(melhor localidade), nesse caso talves os acessos realizados por A, mesmo que menores, tem uma localidade pior.

O programa A talves esteja carregando na cache dados, mas apenas usando um desses dados em cada operacao, precisando entao carregar novos dados para entao realizar suas outras operações.

A estrutura de dados de A pode ser menos otimizada que a de B,como nos casos em que fazer busca de acessos em uma lista gera um custo maior do que fazer em um set.

O programa A pode estar gerando um overhead ao chamar funções que o programa B nao chama

\# FALTA EXEMPLO

\subsection*{Questão 2 -- Hierarquia de memória}
a) Porque o acesso a RAM é mais demorado que o acesso aos caches Ls, que por sua vez são mais demorados que os acessos aos registradores.

b) Um cache miss ocorre quando o dado buscado nao esta no nivel acessado, forçando a CPU a buscar o dado em um nivel abaixo. Eg: buscar em L3 um dado que esta na RAM

c) Um cache hit ocorre quando o dado buscado esta no nivel consultado

d) Localidade Espacial: Os vizinhos dos dados acessados tem a chance de serem acessados.

Localidade Temporal: Os dados agora acessados tem a chance de serem acessados novamente em breve

e) O primeiro padrao de acesso faz uso sequencial, oq gera uma localidade espacial melhor, pois os dados estao carregados em sequencia e tabem uma localidade temporal melhor, pois como o acesso esta ocorrendo sequencialmente, o proximo dado sera acessado em seguida.

O segundo padrao tem um stride muito grande e por isso a memoria cache precisa ser constantemente atualizada, levando a muito mais cache miss

% PART II - Cell 1 (markdown)
\vspace{0.5em}
\section*{Parte II -- Otimização sequencial}

\subsection*{Questão 3 -- Análise do código}

A principal diferença está na alocação de memória: em vez de declarar uma lista vazia e fazer a alocação prévia do espaço de memória, deve-se pré-alocar com tamanho conhecido.

O comando \texttt{append()} realoca toda a memória anterior mais o espaço para o append em cada execução, como o tamanho e o tipo de dado sao conhecidos, fazer a alocação previa é melhor.\\

Fazer 2*math.sqrt(x[i]) para diminuir o overhead.

% PART II - Cell 2+3+4: code cells with functions
\begin{verbatim}
import math 
import numpy as np
from numba import njit

# Versao A 
@njit
def calcula_A(x):
    resultado = []

    for i in range(len(x)):
        valor = math.sqrt(x[i]) + math.sqrt(x[i])
        resultado.append(valor)

    return resultado

# Versao B 
@njit
def calcula_B(x):
    resultado = [0.0] * len(x)
    size = len(x)

    for i in range(size):
        valor = 2 * math.sqrt(x[i])
        resultado[i] = valor

    return resultado

# Versao C
@njit
def calcula_C(x):
    size = len(x)
    resultado = np.empty(size, dtype=np.float64)

    for i in range(size):
        valor = 2 * math.sqrt(x[i])
        resultado[i] = valor

    return resultado
